% ===========================================================================
%  preamble.tex --- packages and macros for the chapter on the
%  birth--death--catastrophe process: distribution theory, quasi-stationarity
%  and burst statistics.
%  When this chapter is dropped into the thesis, merge the contents of this
%  file into the thesis preamble; duplicates may simply be deleted.
% ===========================================================================

\usepackage[T1]{fontenc}
\usepackage[utf8]{inputenc}
\usepackage{lmodern}
\usepackage{microtype}
\usepackage{amsmath,amssymb,amsthm,mathtools}
\usepackage{graphicx}
\usepackage{float}
\usepackage{placeins}
\usepackage{booktabs}
\usepackage{longtable}
\usepackage[hypcap=false]{caption}
\usepackage{enumitem}
\usepackage{xcolor}
\usepackage{array}
\usepackage{hyperref}
\usepackage[nameinlink,noabbrev]{cleveref}

% --- Cross-document references into the mathematical introduction -----------
% The m:... labels live in CH2.  `xr` resolves them against a copy of CH2's
% .aux, so the standalone compile prints real cross-references instead of ??.
% Regenerate ch2.aux whenever CH2 changes:
%     cd "../Good examples/CH2" && latexmk -pdf main.tex && cp main.aux ../../CH5/ch2.aux
% At thesis integration, delete ch2.aux and these two lines: the references
% then resolve natively.
\usepackage{xr}
\externaldocument{ch2}

\graphicspath{{figures/}}

% --- Equation cross-references as bare parenthesised numbers -----------------
\crefformat{equation}{(#2#1#3)}
\Crefformat{equation}{(#2#1#3)}
\crefrangeformat{equation}{(#3#1#4)--(#5#2#6)}
\crefmultiformat{equation}{(#2#1#3)}{ and (#2#1#3)}{, (#2#1#3)}{ and (#2#1#3)}

% --- Notation ---------------------------------------------------------------
% Scalars are set in italic; calligraphic letters carry the release-side
% quantities, so that the load and what it releases are never confused.
\newcommand{\pgf}{probability generating function}
\newcommand{\rv}{random variable}
\newcommand{\pr}[1]{\mathbb{P}\!\left\{#1\right\}}
\newcommand{\E}[1]{\mathbb{E}\!\left[#1\right]}
\newcommand{\var}[1]{\operatorname{Var}\!\left(#1\right)}
\newcommand{\dd}{\mathrm{d}}
\newcommand{\ee}{\mathrm{e}}
\newcommand{\ddt}[1]{\frac{\dd #1}{\dd t}}
\newcommand{\dda}[2]{\frac{\dd #1}{\dd #2}}
\newcommand{\pddt}[1]{\frac{\partial #1}{\partial t}}
\newcommand{\pdd}[2]{\frac{\partial #1}{\partial #2}}
\newcommand{\lrb}[1]{\left(#1\right)}
\newcommand{\lrbs}[1]{\left[#1\right]}
\newcommand{\ind}{\mathbf 1}

% --- Chapter-specific notation regime ---------------------------------------
% Intracellular: birth \lambda, death \mu, catastrophe \delta per particle;
% absorbing rupture state H; load X_t; released count W_t; burst size \Kb
% (never bare K, which is E[X_t^2]); non-fixation probability \Ifix.
\newcommand{\Ifix}[1][]{I_{\mathrm{fix}#1}}
\newcommand{\Kb}{\mathcal{K}}
\newcommand{\Icell}{\mathcal{I}}
\newcommand{\Lap}[1]{\widetilde{#1}}
\newcommand{\Fhyp}{{}_2F_1}
\newcommand{\Xt}{X_t}
\newcommand{\Wt}{W_t}
\newcommand{\ft}{\textit{F.~tularensis}}
\newcommand{\yp}{\textit{Y.~pestis}}

% --- Cross-chapter references -----------------------------------------------
% Redefining these macros is all that thesis integration requires: they become
% "Chapter~\ref{...}" once the chapter order is fixed.  No chapter number
% appears in body prose.
\newcommand{\ChM}{the mathematical introduction}
\newcommand{\ChMcap}{The mathematical introduction}
\newcommand{\ChCore}{the chapter on the birth--death--cata\-strophe process}
\newcommand{\ChCorecap}{The chapter on the birth--death--cata\-strophe process}
\newcommand{\ChPop}{the chapter on population dynamics}
\newcommand{\ChPopcap}{The chapter on population dynamics}
\newcommand{\ChTwoType}{the two-type chapter}
\newcommand{\ChTwoTypecap}{The two-type chapter}
\newcommand{\ChPath}{the chapter on bursting and budding pathogenesis}
\newcommand{\ChPathcap}{The chapter on bursting and budding pathogenesis}

% --- Forward references to later modelling chapters -------------------------
% Every forward reference is routed through \fwd, whose first argument is a
% stable key and whose second is the prose actually typeset.
\newcommand{\fwd}[2]{#2}

% Standalone fallback for labels defined in the mathematical introduction.
% In the assembled thesis the live \cref is used; here the prose fallback
% prints if xr cannot see ch2.aux, so the chapter never ships with ??.
% Keep xr + ch2.aux above: when those labels resolve, \mextref uses \cref.
\providecommand{\mextref}[2]{%
  \ifcsname r@#1\endcsname\cref{#1}\else#2\fi}

% --- Theorem environments ---------------------------------------------------
\theoremstyle{plain}
\newtheorem{theorem}{Theorem}[chapter]
\newtheorem{proposition}[theorem]{Proposition}
\newtheorem{corollary}[theorem]{Corollary}
\theoremstyle{definition}
\newtheorem{definition}[theorem]{Definition}
\theoremstyle{remark}
\newtheorem{remark}[theorem]{Remark}

\allowdisplaybreaks
